\begin{textsample}{NEG Dim 3 – pi – Score 1.00 – pi/pi\_em\_31.txt}  \label{ex:f3_neg_033}
Proposta de aprendizagem em Arte De acordo com o que preconiza a BNCC, é importante consolidar e aprofundar as aprendizagens construídas nas etapas anteriores do Ensino Infantil e do Ensino Fundamental, permitindo ampliar as diferentes formas de comunicação e expressão através da Arte e de diferentes linguagens e saberes, propiciando ao estudante a compreensão das habilidades e competências a serem desenvolvidas no Ensino Médio.
Com a reforma, o currículo será flexibilizado de acordo com os interesses do aluno, sendo que 60 \% da carga horária será de Formação Geral Básica ( BNCC ), obrigatória para todos os estudantes do Ensino Médio ( que irá incorporar o ensino de Arte ) e 40 \% será composta de Itinerários Formativos – em que será ofertado o Componente Arte na área de Linguagens para os alunos que optarem por aprofundar e ampliar seus conhecimentos nesta área.
Para Ana Mae Barbosa ( 2008 ), a reflexão acerca da história é algo primordial para o desenvolvimento de políticas que possam dar suporte para projetar o futuro do ensino-aprendizagem da arte que se deseja ter, mesmo que este não seja um processo imediato.
A Proposta ou Abordagem Triangular que consiste em três dimensões para se construir conhecimentos em arte : contextualização histórica ( conhecer a sua contextualização histórica ) ; fazer artístico ( fazer arte ) ; apreciação artística ( saber ler uma obra de arte ), será concebida na implementação da Base Nacional Comum Curricular, dando seguimento ao que já estava sendo ofertado pelos Parâmetros Curriculares Nacionais ( PCN ).
A abordagem triangular tem significativas apropriações de educadores de outras áreas do conhecimento, pois abre caminhos para o professor na sua prática docente.
Não é um modelo pronto, é flexível, também é permitido mudanças e adequações.
Qualquer educador pode fazer suas escolhas metodológicas, pois não é conteúdo.
A aprendizagem da Arte não pode ser vista apenas como códigos e técnicas, que era concebida no ensino da Arte em décadas passadas.
A BNCC apresenta a proposta de que o produto é tão importante quanto o processo, ambos estão no mesmo patamar.
Nesse contexto, a articulação para a aprendizagem em Arte requer que a abordagem das linguagens articule seis dimensões do conhecimento, que são : criação, crítica, estesia, expressão, fruição e reflexão.
Tais dimensões perpassam os conhecimentos das Artes Visuais, da Dança, da Música, do Teatro e das Artes Integradas, uma novidade na BNCC, que surgiu para articular a interdisciplinaridade entre as linguagens artísticas, possibilitando o uso de novas tecnologias da informação e comunicação.
Desse modo, para que o aluno possa contribuir nesse processo de ensino e aprendizagem é necessário conhecer os parâmetros de suas habilidades, ressignificando a compreensão do fenômeno artístico e suas relações com o mundo contemporâneo a partir de vivências cotidianas, na realização de processos criativos autorais nas linguagens das Artes Visuais, do Audiovisual, da Dança, do Teatro, das Artes Circenses, da Música e das Artes Integradas, considerando as novas tecnologias, como internet e multimídia e seus espaços de compartilhamento e de convívio, estabelecendo relação entre mídia, política, mercado e consumo.
Nessa perspectiva, é importante garantir o respeito às diversas culturas presentes na formação da sociedade brasileira, em especial, as matrizes indígenas e africanas.
Muito se tem discutido, recentemente, acerca da importância de valorizar as expressões da linguagem estética na Arte, contribuindo para o respeito, à diversidade, a sustentabilidade, assegurando uma ação docente efetiva que promova aprendizagens significativas, conhecendo e sabendo utilizar os diferentes procedimentos da produção artística, que envolvam a sensibilidade, a intuição, a subjetividade, o aspecto emocional, a educação inclusiva, o pensamento, as experiências interclasse e extraclasse, para compreender e fazer uma análise crítica acerca da realidade do mundo que o capacita no exercício da cidadania, promovendo um diálogo intercultural, assegurando o seu desenvolvimento pleno das habilidades, sendo capazes da transformação do seu meio através da Arte, com vistas a impulsionar o desenvolvimento integral do aluno.
Nesse contexto, na etapa do Ensino Médio, o professor deve ser o mediador que possibilita o desenvolvimento das competências, levando o aluno a ter um pensamento crítico, argumentativo, investigativo e criativo, a desenvolver seu projeto de vida que o capacita para o mundo do trabalho e da cidadania, a valorizar seu repertório cultural, a conhecer da diversidade, a se conectar ao mundo tecnológico, a ter empatia, a ser colaborativo e responsável, contribuindo na sua formação integral, enfim, a ser agente e protagonista do seu próprio processo de ensino e aprendizagem.

% matched lemmas: 
\end{textsample}
